\documentclass[10pt,a4paper]{amsart}
\usepackage[margin=19mm]{geometry}
\usepackage{amsmath,amssymb,mathtools}
\usepackage[scr=rsfs,cal=euler]{mathalfa}
\usepackage{xfrac}
\usepackage[unicode,hidelinks,bookmarks=false]{hyperref}
\newcommand{\ie}{i.e.\ }
\newcommand{\cf}{cf.\ }
\newcommand{\resp}{respectively\ }
\newcommand{\Subsection}{\subsection}
\providecommand{\nobreakdash}{\nobreak\hbox{-}}
\setlength{\emergencystretch}{2em}
\multlinegap0pt
\usepackage{amssymb}
\usepackage{enumerate}
\usepackage{natbib}
\newtheorem{proposition}{Proposition}
\newtheorem{lemma}{Lemma}
\title{Note on unique representation bases}
\author{Yuchen Ding, Jie Wang}
\date{Source: arXiv:2602.07743v1 (CC BY 4.0)}
\begin{document}
\maketitle

\noindent\emph{Source and licence.} Note on unique representation bases. Version arXiv:2602.07743v1 (CC BY 4.0), distributed by arXiv under the Creative Commons Attribution 4.0 International licence, http://creativecommons.org/licenses/by/4.0/, as recorded in the arXiv licence metadata for this exact version and shown on the abstract page https://arxiv.org/abs/2602.07743v1. Full licence notice: \texttt{licenses/arxiv-2602.07743-CC-BY-4.0.txt}.\par\smallskip

\noindent\textit{Editorial scope.} Counted unit: the article's lemma lem1 (source0.tex lines 207--213). The article's complete original proof of it is delivered with it (source0.tex lines 214--262, one \texttt{proof} environment of 49 lines). No mathematics is written, repaired or re-derived: the statement and the proof are the article's own lines, in the article's own notation, and no retained line is substituted or re-flowed.  Retained uncounted as the source-local context the proof needs: lines 192--194; lines 197--203.  Nothing was omitted from any retained span, so the package carries no omission record.  No retained line contains a citation, so the wrapper carries no bibliography and no bibliography entry.  No retained line contains a figure environment, an image-inclusion command or drawing code, so no image is delivered and the package declares no asset dependency.  Editorial formatting only, in addition: an AMS \texttt{amsart} preamble that restates the article's own declarations and macros used by the retained lines, namely the environments \texttt{proposition}, \texttt{lemma} on separate counters, together with the standard packages those lines need.  The retained environments print in this document as Proposition 1, Lemma 1 in order of appearance, whereas the article prints the same items with the numbering of its own full document.  The complete native source of the article as distributed in the arXiv e-print for this exact version is bundled unchanged as \texttt{sources/194208/source0.tex}.
\begin{proposition}\label{prop1}
As $N\rightarrow\infty$, we have $f_2(N)\sim \sqrt{N}$.
\end{proposition}

\begin{proposition}\label{prop2}
Let $p$ be a prime. Then, there are at least $p$ elements $\{g_1,g_2,\cdots, g_p\}$ in $\big\{1,2,\cdots,p^2-1\big\}$ such that all the sums
$$
g_i+g_j \quad (1\le i\le j\le p)
$$ 
are different modulo $p^2-1$.
\end{proposition}

\begin{lemma}\label{lem1}
Let $\varepsilon>0$ be an arbitrarily small given number. Then, for any sufficiently large prime $p$ there is a Sidon set 
$$
\mathscr{S}\subset \bigg(\left[-p^2, \left(-\frac{1}{2}-\frac{\varepsilon^2}{8}\right)p^2\right]\bigcup \left[\left(\frac{1}{2}+\frac{\varepsilon^2}{8}\right)p^2, p^2\right]\bigg)
$$ 
such that $\big|\mathscr{S}\big|\ge \left(1-\frac{3\varepsilon}{4}\right)p$.
\end{lemma}

\begin{proof}
Let $p$ be a sufficiently large prime. By Proposition \ref{prop2} there are $p$ integers
$$
s_1<\cdots<s_\ell\le \left(\frac{1}{2}-\frac{\varepsilon^2}{8}\right)p^2<s_{\ell+1}<\cdots<s_y<\left(\frac{1}{2}+\frac{\varepsilon^2}{8}\right)p^2\le s_{y+1}<\cdots<s_p
$$
in the interval $\big[1, p^2\big]$ such that 
all sums 
$
s_i+s_j\ (i\le j)
$
are different modulo $p^2-1$. Let
$$
\mathscr{S}=\Big\{s_1-(p^2-1),s_2-(p^2-1),\cdots, s_{\ell-1}-(p^2-1),s_{y+1}, s_{y+2},\cdots, s_p\Big\}.
$$
Then, clearly we have
$$
\mathscr{S}\subset \bigg(\left[-p^2, \left(-\frac{1}{2}-\frac{\varepsilon^2}{8}\right)p^2\right]\bigcup \left[\left(\frac{1}{2}+\frac{\varepsilon^2}{8}\right)p^2, p^2\right]\bigg)
$$ 
Note that $\big\{s_{\ell+1}, s_{\ell+2}, \cdots, s_y\big\}$ is a Sidon set in the interval 
$$
\left[\left(\frac{1}{2}-\frac{\varepsilon^2}{8}\right)p^2, \left(\frac{1}{2}+\frac{\varepsilon^2}{8}\right)p^2\right],
$$ 
which is of length $\frac{\varepsilon^2}{4}p^2$. Thus, by Proposition \ref{prop1} we have
\begin{align}\label{eq-lem-2}
y-\ell\le \big(1+o(1)\big)\sqrt{\frac{\varepsilon^2}{4}p^2}\le \frac{5\varepsilon p}{8}.
\end{align}
Hence, by (\ref{eq-lem-2}) we have
$$
\big|\mathscr{S}\big|=p-(y-\ell)-1\ge p-\frac{5\varepsilon p}{8}-1\ge \left(1-\frac{3\varepsilon}{4}\right)p.
$$

It remains to show that $\mathscr{S}$ is a Sidon set. Let
\begin{align*}
\widetilde{s_i}=
\begin{cases}
s_i-(p^2-1), & \text{if~}1\le i\le \ell-1,\\
s_i, & \text{if~}y+1\le i\le t.
\end{cases}
\end{align*}
Since 
$
s_i+s_j\ (i\le j)
$
are different modulo $p^2-1$, we know that 
$
\widetilde{s_i}+\widetilde{s_j}\ (i\le j)
$
are also different modulo $p^2-1$, which implies that $\mathscr{S}=\big\{\widetilde{s_i}\big\}$ is a Sidon set.
\end{proof}

\end{document}
